% ===== PRÉAMBULE CANONIQUE — à coller VERBATIM en tête de CHAQUE .tex =====
\documentclass[11pt,a4paper]{article}
\usepackage[utf8]{inputenc}\usepackage[T1]{fontenc}\usepackage{lmodern}
\usepackage[french]{babel}
\usepackage{amsmath,amssymb,mathtools}\usepackage{icomma}
\usepackage{siunitx}\sisetup{locale=FR}  % \qty{}{}, \si{}, \ang{}
\usepackage{xcolor}\usepackage{colortbl}
\usepackage{tikz}\usepackage{pgfplots}\pgfplotsset{compat=1.18}
\usepackage[siunitx,europeanresistors]{circuitikz} % circuits (élec.)
\usepackage[most]{tcolorbox}\usepackage{enumitem}\usepackage{array,booktabs}
\usepackage[a4paper,margin=2.2cm]{geometry}
\usetikzlibrary{arrows.meta,calc,angles,quotes,positioning,patterns,%
  backgrounds,shapes.geometric,decorations.pathmorphing,decorations.markings,intersections}
\definecolor{primary}{RGB}{222,49,99}\definecolor{primarydark}{RGB}{150,22,55}
\definecolor{defcol}{RGB}{0,114,178}\definecolor{methcol}{RGB}{52,92,148}
\definecolor{excol}{RGB}{0,135,90}\definecolor{attncol}{RGB}{200,40,55}
\tcbset{boxbase/.style={enhanced,breakable,boxrule=0.9pt,arc=2.5pt,
  left=3mm,right=3mm,top=2.3mm,bottom=2.3mm,fonttitle=\bfseries,coltitle=white}}
% --- boîtes TUTORIEL (noms EXACTS, ne pas renommer) ---
\newtcolorbox{cle}[1][]{boxbase,colback=defcol!6,colframe=defcol,
  title={\textsc{À retenir}\ifx\relax#1\relax\relax\else\textmd{ : \itshape #1}\fi}}
\newtcolorbox{methode}[1][]{boxbase,colback=methcol!7,colframe=methcol,sharp corners,
  title={\textsc{Méthode}\ifx\relax#1\relax\relax\else\textmd{ --- \itshape #1}\fi}}
\newtcolorbox{exemplebox}[1][]{boxbase,colback=excol!5,colframe=excol,
  title={\textsc{Exemple}\ifx\relax#1\relax\relax\else\textmd{ : \itshape #1}\fi}}
\newtcolorbox{piege}[1][]{boxbase,colback=attncol!6,colframe=attncol,
  title={\textsc{Piège}\ifx\relax#1\relax\relax\else\textmd{ : \itshape #1}\fi}}
\newtcolorbox{remarque}[1][]{boxbase,colback=black!4,colframe=black!45,
  title={\textsc{Remarque}\ifx\relax#1\relax\relax\else\textmd{ : \itshape #1}\fi}}
% --- boîtes EXERCICE / CORRIGÉ (noms EXACTS, auto-comptées) ---
\newtcolorbox[auto counter]{exo}[1][]{boxbase,colback=primary!4,colframe=primary,
  title={\textsc{Exercice}~\thetcbcounter\ifx\relax#1\relax\relax\else\textmd{ --- \itshape #1}\fi}}
\newtcolorbox[auto counter]{corrige}[1][]{boxbase,colback=excol!4,colframe=excol,
  title={\textsc{Corrigé}~\thetcbcounter\ifx\relax#1\relax\relax\else\textmd{ --- \itshape #1}\fi}}
\newcommand{\vv}[1]{\vec{#1}}\newcommand{\ds}{\displaystyle}
\newcommand{\R}{\mathbb{R}}\newcommand{\N}{\mathbb{N}}\newcommand{\Z}{\mathbb{Z}}
% (un \usepackage{fancyhdr}+en-tête est possible ; il est ignoré côté web)
% =========================================================================
\begin{document}
\begin{corrige}[deux fils symétriques]
\begin{enumerate}[label=\alph*)]
  \item $G = m\,g = 8\times 10 = \qty{80}{N}$.
  \item L'objet est en équilibre : les fils supportent \emph{tout} le poids, donc la somme des
        contributions verticales des tensions vaut $G=\qty{80}{N}$.
  \item Les deux fils étant identiques et symétriques, ils se partagent la charge à parts égales :
        \[ T = \frac{G}{2} = \frac{80}{2} = \qty{40}{N}. \]
\end{enumerate}
\end{corrige}

\begin{corrige}[une force de plus sur la table]
\begin{enumerate}[label=\alph*)]
  \item Trois forces verticales : le poids $G=m g=4\times 10=\qty{40}{N}$ (bas), la force appliquée
        $F=\qty{12}{N}$ (bas), la réaction $\vv R$ de la table (haut).
  \item Équilibre vertical : $R = G + F = 40 + 12 = \qty{52}{N}$. En appuyant, on augmente la
        réaction du support d'autant.
\end{enumerate}
\end{corrige}

\begin{corrige}[trois forces alignées]
Pour l'équilibre : $\vv F_1+\vv F_2+\vv F_3=\vv 0$. Sur l'axe horizontal (positif vers la droite),
$\vv F_1$ et $\vv F_2$ donnent une résultante $F_2-F_1 = 12-7 = \qty{5}{N}$ vers la droite.
\begin{enumerate}[label=\alph*)]
  \item $\vv F_3$ doit annuler cette résultante : $F_3 = \qty{5}{N}$.
  \item Elle pointe \textbf{vers la gauche} (sens opposé à la résultante de $\vv F_1$ et $\vv F_2$).
\end{enumerate}
\end{corrige}

\begin{corrige}[plaque suspendue : conservation de la charge]
\begin{enumerate}[label=\alph*)]
  \item La charge totale est la somme des indications initiales :
        $G_{\text{tot}} = D_1 + D_2 = 300 + 300 = \qty{600}{N}$.
  \item Le poids total ne change pas ; à l'équilibre, $D_1+D_2$ reste égal à \qty{600}{N}. Donc
        \[ D_2 = 600 - D_1 = 600 - 220 = \qty{380}{N}. \]
        Ce que $D_1$ perd ($\qty{80}{N}$), $D_2$ le gagne : leur somme est conservée.
\end{enumerate}
\end{corrige}

\begin{corrige}[deux fils asymétriques]
\begin{enumerate}[label=\alph*)]
  \item \textbf{Non.} Quelle que soit l'inclinaison, l'objet reste suspendu, donc son poids est
        toujours intégralement porté : la somme des contributions verticales vaut $G$, constant.
  \item Le \textbf{fil de gauche} (le plus proche de la verticale) affiche la plus grande valeur. Un fil
        presque vertical fournit une contribution verticale importante avec peu d'inclinaison ; il
        supporte donc la plus grande part du poids. Le fil très incliné en supporte moins.
\end{enumerate}
\end{corrige}

\end{document}
